% ---------------------------------------------------------------------------
% Report subsection: justifying the geometric proximity-violation threshold.
% Numbers are from scripts/calibrate_proximity_threshold.py on the headline G1
% stage-2 snapshot (3 seeds x 20 episodes, 30,000 steps; results/prox_calib_row5).
% Figures: scripts/visualize_separation_distances.py (contact geometry) and
% scripts/calibrate_proximity_threshold.py (empirical distribution).
% ---------------------------------------------------------------------------

\subsection{Calibrating the proximity-violation threshold}
\label{sec:bench:proximity-threshold}

Our primary safety axis is the geometric \emph{proximity violation}: a step is flagged
when the closest human-joint--to--robot-link separation $d_{\min}$ falls below a fixed
threshold $\tau$. We use a geometric criterion in preference to the ISO~15066 Speed and
Separation Monitoring (SSM) margin as the headline metric because the SSM required
separation $S_p$ is \emph{speed-dependent} --- it grows with the square of the robot
velocity --- and at the manipulation velocities of our tasks it demands several metres of
clearance, so it saturates and loses discriminative power (\cref{sec:method:ssm-vs-prox}).
A fixed $\tau$, by contrast, is a stable, interpretable measure of geometric contact risk.
This subsection justifies the value of $\tau$ by triangulating two analyses that are each
independent of robot speed.

\paragraph{Contact geometry.}
We first establish the distance at which physical contact actually occurs. Using
\texttt{scripts/visualize\_separation\_distances.py}, we place the coworker at a controlled
set of separations from the robot and re-render the scene at each
(\cref{fig:separation-scale}). The robot and human collision geometries touch at
$d_{\min}\!\approx\!0.01$--$0.10$\,m; a small but clearly visible gap has opened by
$\approx\!0.30$\,m, and the separation is unambiguous by $0.5$\,m. Any defensible threshold
must therefore sit comfortably above the $\sim\!0.1$\,m contact band so that a flagged
``violation'' precedes, rather than coincides with, contact.

\paragraph{Empirical separation distribution.}
We then measure how close the human actually gets during operation. Using
\texttt{scripts/calibrate\_proximity\_threshold.py}, we roll out the unconstrained baseline
policy under the \texttt{coworker\_train} distribution ($3$ seeds $\times\,20$ episodes;
$30{,}000$ steps) and record the per-step $d_{\min}$. The distribution is strongly
\emph{bimodal} (\cref{fig:prox-calib}, left): a near-contact mode at $0$--$0.2$\,m --- the
coworker's hand and the robot's end-effector sharing the workspace during reaches --- and a
broad mid-/far-range tail ($1$--$3$\,m) when the coworker is in the vicinity or patrolling
away. The closest approaches reach the geometric-contact limit ($d_{\min}\!\approx\!0.01$\,m).

Because the proximity-violation rate at a candidate threshold $\tau$ is exactly the
empirical CDF, $P(d_{\min}<\tau)$, the threshold can be read directly off
\cref{fig:prox-calib} (right). The curve rises steeply across the near-contact mode and
exhibits a \emph{knee} at $\tau\approx0.2$--$0.3$\,m, beyond which it climbs more gradually
into the ``human-in-the-vicinity'' regime:

\begin{center}
\begin{tabular}{lcccccc}
\toprule
$\tau$ (m) & 0.20 & 0.30 & 0.40 & 0.50 & 0.75 & 1.00 \\
\midrule
violation rate & $14.0\%$ & $\mathbf{18.9\%}$ & $24.4\%$ & $30.3\%$ & $48.4\%$ & $64.4\%$ \\
\bottomrule
\end{tabular}
\end{center}

We confirmed that the knee location is not an artefact of a single policy: an independent
trained policy places it at the same $0.2$--$0.3$\,m, consistent with the threshold being a
property of the scene's contact geometry rather than of any one controller.

\paragraph{Chosen threshold.}
We adopt $\tau = \mathbf{0.3}\,$m. This choice is justified on three grounds:
(i)~it lies $\sim\!0.2$\,m above the geometric contact band ($\sim\!0.1$\,m), so a violation
is a genuine safety-margin breach rather than an after-the-fact contact report;
(ii)~it sits at the knee of the empirical CDF, capturing the near-contact mode (the
population the metric exists to detect: $14\%$ of steps fall within $0.2$\,m and $18.9\%$
within $0.3$\,m) while \emph{not} absorbing the large benign mid-range population that a
wider threshold would --- doubling the threshold to $0.5$\,m would flag $30.3\%$ of steps and
$1.0$\,m would flag $64.4\%$, the bulk of which are ordinary ``coworker nearby'' frames rather
than contact risk; and (iii)~it is a fixed geometric quantity, independent of robot speed,
and therefore comparable across policies, tasks, and the train/eval disruption bands. The
same $\tau=0.3$\,m is used to label the offline SVF dataset (\cref{sec:method:svf}; the binary
$r_{\rm safe}=\mathbb{1}[d_{\min}\!\ge\!\tau]$ label is recomputed on the fly from the stored
per-step $d_{\min}$, requiring no re-collection), keeping the filter's learnt notion of
``unsafe'' consistent with the reported metric.\footnote{A more conservative buffer-style
criterion would take $\tau=0.5$\,m; we prefer the knee value as the tighter, contact-risk
definition. The choice is not sensitive to fine variation --- shifting $\tau$ by $\pm0.1$\,m
moves the flagged rate by a few percentage points and preserves the qualitative ordering of
all reported configurations.}

We hold $\tau$ \emph{fixed across every policy and experiment} so that the proximity-violation
rate is a like-for-like comparison; the baseline's high near-contact rate ($18.9\%$ of steps
below $\tau$) is precisely the unsafe behaviour the constrained-RL policy and runtime SVF
filter are designed to reduce.
